% Irodov Part V §5.3 originals 119--130.
% Uploaded source pp.216--218 (printed 208--209); source answers pp.330--331.
\diffractionproblem{119}{208}
波长为 $\lambda$ 的平面单色光垂直入射长矩形单缝，缝宽为 $b$。
求夫琅禾费远场光强随衍射角 $\theta$ 的分布与暗纹位置。
\exsolution
取缝宽横坐标 $x\in[-b/2,b/2]$。
方向 $\theta$ 上的复振幅为均匀孔径相量积分
\[
\widetilde E(\theta)\propto
\int_{-b/2}^{b/2}e^{ikx\sin\theta}\,\mathrm dx
=b\frac{\sin\alpha}{\alpha},
\qquad
\alpha=\frac{\pi b}{\lambda}\sin\theta.
\]
平方得
\[
\boxed{I(\theta)=I(0)
\left(\frac{\sin\alpha}{\alpha}\right)^2}.
\]
除中央主极大以外，
暗纹要求 $\sin\alpha=0$ 且 $\alpha\ne0$，
\[
\boxed{b\sin\theta=m\lambda},
\quad m=\pm1,\pm2,\ldots .
\]
\excheck $\alpha\to0$ 时
$(\sin\alpha/\alpha)^2\to1$，
中央主极大不应因 $\sin0=0$ 而误判为暗纹。

\diffractionproblem{120}{208}
利用第119题结果求单缝夫琅禾费衍射的第一、第二、第三\emph{副极大}角位置。
\exsolution
令
$f(\alpha)=(\sin\alpha/\alpha)^2$，
对非零且非暗纹处求导：
$f'(\alpha)\propto
\sin\alpha(\alpha\cos\alpha-\sin\alpha)$。
因此副极大必须满足
\[
\boxed{\tan\alpha=\alpha}.
\]
逐区间解此超越方程的前三个正根为
\[
\alpha_1\simeq1.4303\pi,\quad
\alpha_2\simeq2.4590\pi,\quad
\alpha_3\simeq3.4709\pi.
\]
因 $\alpha=\pi b\sin\theta/\lambda$，
\[
\boxed{b\sin\theta_1\simeq1.43\lambda},\quad
\boxed{b\sin\theta_2\simeq2.46\lambda},\quad
\boxed{b\sin\theta_3\simeq3.47\lambda}.
\]
\excheck $\alpha=(m+1/2)\pi$ 只是近似，不是副极大严格条件；高阶根趋近此近似。

\diffractionproblem{121}{208}
波长
$\lambda=0.50\,\mathrm{\mu m}$ 的光以
$\theta_0=30^\circ$ 入射宽
$b=10\,\mathrm{\mu m}$ 的单缝。
求中央夫琅禾费极大两侧相邻第一暗纹的衍射角。
\exsolution
斜入射的相位差由
$k(\sin\theta-\sin\theta_0)x$ 给出，
故单缝第一暗纹满足
\[
\boxed{b(\sin\theta-\sin\theta_0)=\pm\lambda}.
\]
代
$\sin\theta_0=0.50$、
$\lambda/b=0.050$，
\[
\theta_-=\arcsin(0.450)\simeq26.7^\circ,\quad
\theta_+=\arcsin(0.550)\simeq33.4^\circ .
\]
因此
\[
\boxed{\theta_-\simeq27^\circ,\qquad
\theta_+\simeq33^\circ}.
\]
\excheck 它们关于 $\theta_0$ 的角度并非精确对称，
因为 $\sin\theta$ 是非线性函数。

\diffractionproblem{122}{208}
波长 $\lambda=0.60\,\mathrm{\mu m}$ 的平面光垂直照到玻璃楔的第一表面，
楔角 $\theta=15^\circ$；
相反面不透明但有一条宽
$b=10\,\mathrm{\mu m}$、平行楔棱的单缝。
求：(a) 零级衍射主极大相对原入射光的偏转角；
(b) 零级主极大的角宽。
\exsolution
假定玻璃 $n\simeq1.50$。
楔内光在第二表面的入射角为 $\theta$，
折出空气的中央方向满足
$\sin\beta=n\sin\theta$。
(a) 相对原光线的偏折量
\[
\boxed{\Delta\theta_0=\arcsin(n\sin\theta)-\theta
\simeq7.9^\circ}.
\]
(b) 相对内部相位中心，屏外第一暗纹满足
$b[\sin\beta_\pm-n\sin\theta]=\pm\lambda$，
所以零级角宽
\[
\boxed{\Delta\theta_{\rm width}=
\arcsin\left(n\sin\theta+\frac{\lambda}{b}\right)
-\arcsin\left(n\sin\theta-\frac{\lambda}{b}\right)
\simeq7.3^\circ}.
\]
\excheck (a) 是中央方向\emph{偏转量}，
(b) 是中央主瓣两侧第一暗纹的\emph{全角宽}，两者物理含义不同。

\diffractionproblem{123}{208}
平面单色光以掠射角
$\alpha_0=1.0^\circ$ 入射周期
$d=1.0$mm 的反射光栅；
若在掠射角 $\alpha=3.0^\circ$ 方向出现二级夫琅禾费主极大，求波长。
\exsolution
反射式光栅的掠射角相位差为
$d(\cos\alpha_0-\cos\alpha)$，
二级极大满足
\[
2\lambda=d(\cos\alpha_0-\cos\alpha),
\]
故
\[
\boxed{\lambda=
\frac d2(\cos1^\circ-\cos3^\circ)
\simeq0.61\,\mathrm{\mu m}}.
\]
\excheck 掠射角是光线与\emph{光栅面}的夹角，
不能把它当与法线的入射角直接套正弦公式。

\diffractionproblem{124}{208}
对由三条等宽缝组成的透射光栅，作夫琅禾费衍射的强度分布草图，
其中光栅周期 $d$ 与缝宽 $b$ 之比为
(a) $d/b=2$；
(b) $d/b=3$。
指出缺级和副极大数目。
\exsolution
三缝场为单缝衍射包络乘三等距缝相干和，
\[
\boxed{I(\theta)=I_1
\left(\frac{\sin\alpha}{\alpha}\right)^2
\left(\frac{\sin3\beta}{\sin\beta}\right)^2},
\quad
\alpha=\frac{\pi b\sin\theta}{\lambda},\
\beta=\frac{\pi d\sin\theta}{\lambda}.
\]
主极大在 $d\sin\theta=m\lambda$，
包络零点在 $b\sin\theta=p\lambda$。
(a) 若 $d/b=2$，
主极大的偶数级 $m=2p$ 同时落在包络零点，
$\boxed{\text{偶数级缺级}}$。
(b) 若 $d/b=3$，
$m=3p$ 缺级，
$\boxed{\text{每第三级缺级}}$。
三缝干涉因子在相邻两个主极大之间有
$N-2=1$ 个副极大、$N-1=2$ 个附加极小；
依此绘制包络与主峰即可。
\excheck 只给主峰位置而不乘单缝包络会漏掉缺级现象。

\diffractionproblem{125}{208}
单色光垂直入射光栅，
$\lambda_1=0.65\,\mathrm{\mu m}$ 的二级极大出现在
$45^\circ$；
求波长
$\lambda_2=0.50\,\mathrm{\mu m}$ 的三级极大角度。
\exsolution
由
$d\sin45^\circ=2\lambda_1$，
得
$d=2\lambda_1/\sin45^\circ$。
同一光栅的三级满足
\[
\sin\theta_3=\frac{3\lambda_2}{d}
=\frac{3\lambda_2\sin45^\circ}{2\lambda_1}.
\]
所以
\[
\boxed{\theta_3=
\arcsin\frac{3(0.50)\sin45^\circ}{2(0.65)}
\simeq55^\circ}.
\]
\excheck 必须确认 $\sin\theta_3<1$，否则该衍射级不存在。

\diffractionproblem{126}{208}
波长 $\lambda=535$nm 的光垂直入射光栅。
$35^\circ$ 是某一级主极大方向，
且最高可见衍射级数为5。
求光栅周期 $d$。
\exsolution
设 $35^\circ$ 对应级次 $m$，
$d\sin35^\circ=m\lambda$，
所以 $d/\lambda=m/\sin35^\circ$。
最高级数为
$\lfloor d/\lambda\rfloor=5$，
即
$5\le m/\sin35^\circ<6$，
$5\sin35^\circ\le m<6\sin35^\circ$，
只有整数 $m=3$ 可满足。
因此
\[
\boxed{d=\frac{3\lambda}{\sin35^\circ}
\simeq2.8\,\mathrm{\mu m}}.
\]
\excheck 角度为35°的级次\emph{不等于最高级次5}，
应先用最高级次范围反推出 $m=3$。

\diffractionproblem{127}{208}
波长未知的单色光垂直落在周期
$d=2.2\,\mathrm{\mu m}$ 的光栅，
同一侧一级和二级主极大方向的角距为
$\Delta\theta=15^\circ$。
求波长。
\exsolution
设 $\theta_2=\theta_1+\Delta\theta$，
光栅方程
$\sin\theta_1=\lambda/d=x$，
$\sin\theta_2=2x$。
由正弦加角公式
$2x=x\cos\Delta\theta+
\sqrt{1-x^2}\sin\Delta\theta$，
解得
\[
x=\frac{\sin\Delta\theta}
{\sqrt{5-4\cos\Delta\theta}}.
\]
因此
\[
\boxed{\lambda=
\frac{d\sin15^\circ}
{\sqrt{5-4\cos15^\circ}}
\simeq0.54\,\mathrm{\mu m}}.
\]
\excheck 两级必须同侧、且二级出射角小于90°；计算得到的解满足这一条件。

\diffractionproblem{128}{208}
波长 $\lambda=530$nm 的光落到周期
$d=1.50\,\mathrm{\mu m}$ 的透射光栅。
求\emph{最高级}主极大相对光栅法线的方向角，分别在
(a) 垂直入射；
(b) 以 $60^\circ$ 倾角入射时。
\exsolution
透射光栅方程
\[
d(\sin\theta_m-\sin\theta_0)=m\lambda
\]
（有向 $\theta_m$）。
(a) $\theta_0=0$，
最高 $m_{\max}=\lfloor d/\lambda\rfloor=2$，
\[
\boxed{\theta_2=\arcsin(2\lambda/d)
\simeq45^\circ}.
\]
(b) $\theta_0=60^\circ$，
可见角 $\sin\theta\in[-1,1]$；
负向可见级次下限
$m_{\min}=\lceil(-1-\sin60^\circ)d/\lambda\rceil=-5$。
最高绝对级的 $m=-5$ 对应
\[
\boxed{\theta_{-5}=
\arcsin\left(\sin60^\circ-\frac{5\lambda}{d}\right)
\simeq-65^\circ}.
\]
\excheck 斜入射时不同侧允许的最大正、负级次不对称，
不能只用 $\lfloor d/\lambda\rfloor$。

\diffractionproblem{129}{209}
在平凸柱面玻璃透镜的平面上刻透射光栅，
光栅周期 $d=6.0\,\mathrm{\mu m}$，
柱面曲率半径 $R=20$cm。
波长 $\lambda=0.60\,\mathrm{\mu m}$ 的平面光垂直入射，
求焦平面内左右对称的两个一级主极大之间的距离。
\exsolution
普通玻璃取 $n\simeq1.50$，
平凸柱面透镜在聚焦方向焦距
$f=R/(n-1)=0.40$m。
一级衍射角满足
$\sin\theta_1=\lambda/d=0.10$。
焦面左右极大坐标为
$x_\pm=\pm f\tan\theta_1$，故
\[
\boxed{\Delta x=2f\tan(\arcsin(\lambda/d))
\simeq2(0.40)\frac{0.10}{\sqrt{1-0.10^2}}
\simeq8.0\,\mathrm{cm}}.
\]
\excheck 本题是\emph{柱面}透镜，焦线沿光栅刻线延展，但垂直刻线的坐标间距仍可用上述焦面几何求解。

\diffractionproblem{130}{209}
波长 $\lambda=0.50\,\mathrm{\mu m}$ 的平面波垂直入射顶角
$A=30^\circ$ 的玻璃楔前表面；
在另一面刻透射光栅，周期
$d=2.00\,\mathrm{\mu m}$，
刻线平行楔棱。
求零级和一级主极大与原入射光的夹角，
并求最高级主极大的级次及角位置。
\exsolution
玻璃取 $n\simeq1.50$。
原光沿前表面法线，在后表面上的内部入射角为 $A$。
透射光栅相位条件可写为
\[
d[n\sin A-\sin(A-\theta_m)]=m\lambda,
\]
其中 $\theta_m$ 以原入射方向为有向零点。
因此
$\sin(A+\theta_m)=n\sin A-m\lambda/d
=0.75-0.25m$；
其中 $\theta_m$ 是按原书图示从初始入射光向出射光测量的\emph{有向}角。
由
$\theta_m=A-\arcsin(0.75-0.25m)$，
得到
\[
\boxed{\theta_0\simeq-18.5^\circ},\quad
\boxed{\theta_{+1}=0^\circ}.
\]
在原图可出射方向范围内可见最高正级
$\boxed{m_{\max}=+6}$，其对原入射方向的角度
\[
\boxed{\theta_{+6}\simeq+78.5^\circ}.
\]
\excheck 源题角号依图的具体零向量正方向决定，光栅方程应先给出；级次不能仅凭 $d/\lambda=4$ 判定，因为光线在楔内斜入射。
